\documentclass[12pt,a4paper]{article}
\usepackage[UTF8]{ctex}          % 添加中文支持
\usepackage{amsmath,amssymb,amsfonts,geometry,enumitem,graphicx}
\geometry{left=2.5cm,right=2.5cm,top=2.5cm,bottom=2.5cm}
\title{2024年普通高等学校招生全国统一考试\\
       \large 新高考Ⅰ卷·数学}
\date{}
\begin{document}
\maketitle

\begin{center}
{\large 本试卷共4页，19小题，满分150分。考试用时120分钟。}
\end{center}

\section*{一、选择题：本题共8小题，每小题5分，共40分。在每小题给出的四个选项中，只有一项是符合题目要求的。}

\begin{enumerate}[label=\arabic*.]
\item 已知集合 $A=\{x \mid -5 < x^3 < 5\}$，$B=\{-3,-1,0,2,3\}$，则 $A \cap B =$ \hfill 【　】

A. $\{-1,0\}$ \quad B. $\{2,3\}$ \quad C. $\{-3,-1,0\}$ \quad D. $\{-1,0,2\}$

\vspace{0.3cm}

\item 若 $\dfrac{z}{z-1}=1+i$，则 $z=$ \hfill 【　】

A. $-1-i$ \quad B. $-1+i$ \quad C. $1-i$ \quad D. $1+i$

\vspace{0.3cm}

\item 已知向量 $\vec{a}=(0,1)$，$\vec{b}=(2,x)$，若 $\vec{b} \perp (\vec{b}-4\vec{a})$，则 $x=$ \hfill 【　】

A. $-2$ \quad B. $-1$ \quad C. $1$ \quad D. $2$

\vspace{0.3cm}

\item 已知 $\cos(\alpha+\beta)=m$，$\tan\alpha\tan\beta=2$，则 $\cos(\alpha-\beta)=$ \hfill 【　】

A. $-3m$ \quad B. $-\dfrac{m}{3}$ \quad C. $\dfrac{m}{3}$ \quad D. $3m$

\vspace{0.3cm}

\item 已知圆柱和圆锥的底面半径相等，侧面积相等，且它们的高均为$\sqrt{3}$，则圆锥的体积为 \hfill 【　】

A. $2\sqrt{3}\pi$ \quad B. $3\sqrt{3}\pi$ \quad C. $6\sqrt{3}\pi$ \quad D. $9\sqrt{3}\pi$

\vspace{0.3cm}

\item 已知函数 $f(x)=\begin{cases} -x^2-2ax-a, & x<0 \\ e^x+\ln(x+1), & x\geqslant 0 \end{cases}$ 在$\mathbf{R}$上单调递增，则$a$的取值范围是 \hfill 【　】

A. $(-\infty,0]$ \quad B. $[-1,0]$ \quad C. $[-1,1]$ \quad D. $[0,+\infty)$

\vspace{0.3cm}

\item 当 $x\in[0,2\pi]$ 时，曲线 $y=\sin x$ 与 $y=2\sin\left(3x-\dfrac{\pi}{6}\right)$ 的交点个数为 \hfill 【　】

A. 3 \quad B. 4 \quad C. 6 \quad D. 8

\vspace{0.3cm}

\item 已知函数 $f(x)$ 的定义域为 $\mathbf{R}$，$f(x)>f(x-1)+f(x-2)$，且当 $x<3$ 时 $f(x)=x$，则下列结论中一定正确的是 \hfill 【　】

A. $f(10)>100$ \quad B. $f(20)>1000$ \quad C. $f(10)<1000$ \quad D. $f(20)<10000$

\end{enumerate}

\section*{二、选择题：本题共3小题，每小题6分，共18分。在每小题给出的选项中，有多项符合题目要求。全部选对的得6分，部分选对的得部分分，有选错的得0分。}

\begin{enumerate}[label=\arabic*.]
\item 为了解推动出口后的亩收入（单位：万元）情况，从该种植区抽取样本，得到推动出口后亩收入的样本均值为$\overline{x}=2.1$，样本方差$s^2=0.01$。已知该种植区以往的亩收入$X$服从正态分布$N(1.8,0.1^2)$，假设推动出口后的亩收入$Y$服从正态分布$N(\overline{x},s^2)$，则 \hfill 【　】

A. $P(X>2)>0.2$ \quad B. $P(X>2)<0.5$ \\
C. $P(Y>2)>0.5$ \quad D. $P(Y>2)<0.8$

\vspace{0.3cm}

\item 设函数 $f(x)=(x-1)^2(x-4)$，则 \hfill 【　】

A. $x=3$ 是 $f(x)$ 的极小值点 \\
B. 当 $0<x<1$ 时，$f(x)<f(x^2)$ \\
C. 当 $1<x<2$ 时，$-4<f(2x-1)<0$ \\
D. 当 $-1<x<0$ 时，$f(2-x)>f(x)$

\vspace{0.3cm}

\item 造型 $\alpha$ 可以看作图中的曲线 $C$ 的一部分。已知 $C$ 过坐标原点 $O$，且 $C$ 上的点满足到两个定点 $F_1(-a,0)$、$F_2(a,0)$（$a>0$）的距离之积为 $a^2$，则 \hfill 【　】

\begin{quote}
\emph{【图片描述：图中展示了一条形似横放的“8”字或双扭线的封闭曲线，左右对称，中心在原点，曲线在原点处有交叉，整体呈横向伸展的蝴蝶结状。】}
\end{quote}

A. 点 $(0,0)$ 到 $C$ 上点的距离的最小值为 $a$ \\
B. 曲线 $C$ 上的点的横坐标的取值范围是 $[-2a,2a]$ \\
C. 点 $(a,0)$ 到 $C$ 上点的距离的最小值为 $a$ \\
D. 曲线 $C$ 上存在点 $P$，使得 $\triangle PF_1F_2$ 的面积为 $a^2$

\end{enumerate}

\section*{三、填空题：本题共3小题，每小题5分，共15分。}

\begin{enumerate}[label=\arabic*.]
\item 设双曲线 $C:\dfrac{x^2}{a^2}-\dfrac{y^2}{b^2}=1$（$a>0,b>0$）的左、右焦点分别为 $F_1,F_2$，过 $F_2$ 作平行于 $y$ 轴的直线交 $C$ 于 $A,B$ 两点，若 $|F_1A|=13$，$|AB|=10$，则 $C$ 的离心率为 \_\_\_\_\_\_\_\_。

\vspace{0.3cm}

\item 若曲线 $y=e^x+x$ 在点 $(0,1)$ 处的切线也是曲线 $y=\ln(x+1)+a$ 的切线，则 $a=$ \_\_\_\_\_\_\_\_。

\vspace{0.3cm}

\item 甲、乙两人从 $1,2,\cdots,8$ 中各任选一个数，甲先选（不放回），乙后选，则甲选到的数大于乙选到的数的概率为 \_\_\_\_\_\_\_\_。

\end{enumerate}

\section*{四、解答题：本题共5小题，共77分。解答应写出文字说明、证明过程或演算步骤。}

\begin{enumerate}[label=\arabic*.]
\item （13分）记 $\triangle ABC$ 的内角 $A,B,C$ 的对边分别为 $a,b,c$，已知 $\sin C=\sqrt{2}\cos B$，$a^2+b^2-c^2=\sqrt{2}ab$。
\begin{enumerate}[label=(\arabic*)]
\item 求 $B$；
\item 若 $\triangle ABC$ 的面积为 $3+\sqrt{3}$，求 $c$。
\end{enumerate}

\vspace{0.3cm}

\item （15分）已知 $A(0,3)$ 和 $P(3,\frac{3}{2})$ 为椭圆 $C:\dfrac{x^2}{a^2}+\dfrac{y^2}{b^2}=1$（$a>b>0$）上的两点。
\begin{enumerate}[label=(\arabic*)]
\item 求 $C$ 的离心率；
\item 过点 $P$ 的直线 $l$ 交 $C$ 于另一点 $B$，且 $\triangle ABP$ 的面积为 $9$，求 $l$ 的方程。
\end{enumerate}

\vspace{0.3cm}

\item （15分）如图，四棱锥 $P-ABCD$ 中，$PA\perp$ 底面 $ABCD$，$PA=AC=2$，$BC=1$，$AB=\sqrt{3}$。
\begin{quote}
\emph{【图片描述：四棱锥$P-ABCD$，底面$ABCD$为四边形，$PA$垂直于底面，$AC$为底面的一条对角线，$AB$与$BC$在底面上呈一定角度。】}
\end{quote}
\begin{enumerate}[label=(\arabic*)]
\item 若 $AD\perp PB$，证明：$AD\parallel$ 平面 $PBC$；
\item 若 $AD\perp DC$，且二面角 $A-CP-D$ 的正弦值为 $\dfrac{\sqrt{42}}{7}$，求 $AD$。
\end{enumerate}

\vspace{0.3cm}

\item （17分）已知函数 $f(x)=\ln\dfrac{x}{2-x}+ax+b(x-1)^3$。
\begin{enumerate}[label=(\arabic*)]
\item 若 $b=0$，且 $f'(x)\geqslant 0$，求 $a$ 的最小值；
\item 证明：曲线 $y=f(x)$ 是中心对称图形；
\item 若 $f(x)>-2$ 当且仅当 $1<x<2$，求 $b$ 的取值范围。
\end{enumerate}

\vspace{0.3cm}

\item （17分）设 $m$ 为正整数，数列 $a_1,a_2,\cdots,a_{4m+2}$ 是公差不为 $0$ 的等差数列，若从中删去两项 $a_i$ 和 $a_j$（$i<j$）后剩余的 $4m$ 项可被平均分为 $m$ 组，且每组中4个数能构成等比数列，则称数列 $a_1,a_2,\cdots,a_{4m+2}$ 是 $(i,j)$-可分数列。
\begin{enumerate}[label=(\arabic*)]
\item 写出所有的 $(i,j)$，$1\leqslant i<j\leqslant 6$，使数列 $a_1,a_2,\cdots,a_6$ 是 $(i,j)$-可分数列；
\item 当 $m\geqslant 3$ 时，证明：数列 $a_1,a_2,\cdots,a_{4m+2}$ 是 $(2,13)$-可分数列；
\item 从 $1,2,\cdots,4m+2$ 中一次任取两个数 $i$ 和 $j$（$i<j$），记使数列 $a_1,a_2,\cdots,a_{4m+2}$ 是 $(i,j)$-可分数列的概率为 $P_m$，证明：$P_m>\dfrac{1}{8}$。
\end{enumerate}

\end{enumerate}

\end{document}